\section*{The Propers: Detailed Commentary}

The commentary proceeds element by element, in the order the missal prints
them. It opens with the complete appointed formulary, so that every claim made
afterwards can be checked against the text it is made about; then it reports
what two 1962 editions actually print and where the appointed text differs from
the Bible's; then it treats each element in turn, moving from the complete
canonical context of the passage, through the direct exegesis of witnesses who
had it in front of them, to later reception --- explaining what each witness
actually argues and preserving disagreement where the tradition disagrees.

\subsection*{The appointed formulary in full}

Everything the 1962 \latin{Missale Romanum} prints for this Mass, in the order
it prints it: the Latin of the Vatican \latin{editio typica} exactly as the
controlling facsimile has it, each element under its own printed heading,
scriptural reference and marginal number, and beneath it the English of an
identified public-domain witness. This is the canonical full research edition;
a reader who came for the texts has them here.

\textbf{Formulary.} \latin{DOMINICA UNDECIMA post Pentecosten}, \latin{II
classis}, printed pp.~390--392, marginal nos.\ 1542--1551. Green follows the
general rubrics for Sundays after Pentecost and is not printed on those pages.
The formulary begins immediately after the Postcommunion of the Tenth Sunday
(no.\ 1541) and ends immediately before the heading \latin{DOMINICA DUODECIMA
post Pentecosten} (no.\ 1552). Nothing between those boundaries is omitted
here. It appoints no Tract, Sequence, second oration, blessing or ritual text,
and no second Collect, Secret or Postcommunion.

\textbf{The Latin.} Received liturgical text of the Roman Rite, older than the
edition that prints it. No public-domain conclusion is claimed for the 1962
typical edition itself; what permits the Latin here is 17 U.S.C.\ \S103(b),
under which a revised edition's copyright reaches only that edition's own
contribution and not the preexisting material it reprints. All ten appointed
elements of this formulary are printed in the Pustet Missal of Ratisbon, 1862,
and all ten again in the Venice Missal of 1570, both of which this repository
tracks as public-domain text; the terminal appendix states the basis and its
bounds in full. The project claims no rights in this text and has not
composed, translated, modernised, conflated or adapted any of it. Every form
below was visually collated at 300 and 600 dpi against the CMAA facsimile of
the Vatican typical edition, and independently against page images of the 1962
Benziger \latin{editio iuxta typicam}. Where the two editions differ, the
Vatican typical reading is what appears here and the difference is tabulated
in the collation below.

\textbf{The English.} Scriptural elements are the Douay--Rheims as revised by
Bishop Challoner (1749--1752) --- the English of the Clementine Vulgate the
missal prints --- in the Project Gutenberg e-text of that revision. The three
orations are the anonymous English of \work{The Roman Missal translated into
the English language for the use of the laity}, first revised edition
(Philadelphia: Eugene Cummiskey, 1861), from the three lines that book
supplies for this formulary, reproduced with its own nineteenth-century
spelling and its abbreviated conclusion ``Thro'.'' Both are public domain in
the United States. Neither is an official liturgical translation, and neither
is the text of the 1962 books: this is a study companion, not an altar book or
a hand missal. The project has translated nothing. Where the registered
English does not answer the missal's Latin --- and it fails to at several
points in this formulary --- the gap is marked at the element and left open.

\elementguard
\subsubsection*{1. Introit \cue{Int.}}

\begin{appointedlatin}{\latin{Ant. ad Introitum} --- Ps. 67, 6-7 et 36 --- marginal no.\ 1542}
Deus in loco sancto suo: Deus qui inhabitáre facit unánimes in domo: ipse
dabit virtútem et fortitúdinem plebi suæ. \textnormal{Ps. ibid., 2} Exsúrgat
Deus, et dissipéntur inimíci eius: et fúgiant, qui odérunt eum, a fácie eius.
℣.~Glória Patri.
\end{appointedlatin}

\begin{namedtranslation}{English: Douay--Rheims (Challoner), Ps. 67:6--7 and 36 in full, with the fragments the antiphon takes in bold, and Ps. 67:2}
\textbf{v.~6} Who is the father of orphans, and the judge of widows.
\textbf{God in his holy place:} \textbf{v.~7} \textbf{God who maketh men of
one manner to dwell in a house:} Who bringeth out them that were bound in
strength; in like manner them that provoke, that dwell in sepulchres.
\textbf{v.~36} God is wonderful in his saints: the God of Israel is
\textbf{he who will give power and strength to his people.} Blessed be God.

\textit{Ps.}~(v.~2) Let God arise, and let his enemies be scattered: and let
them that hate him flee from before his face. ℣.~\latin{Glória Patri.}
\end{namedtranslation}

\englishgap{Two declarations belong beside this text. \textit{First}, the
antiphon is a centonisation, which is why the missal's own reference reads
\latin{6-7 et 36}: it takes the closing clause of v.~6, the opening clause of
v.~7, and the central clause of v.~36, so the three Douay verses are printed
whole with the sung fragments marked rather than stitched into a seamless
English sentence the project would itself have composed. \textit{Second}, at
v.~7 the antiphon sings \latin{unánimes} --- ``those of one mind'' --- where
the Clementine, and therefore the Douay, reads \latin{unius moris}, ``of one
manner.'' The Douay cannot be quoted for the chant's word, and this guide does
not supply a rendering of its own; the older-psalter reading behind
\latin{unánimes} is discussed in the element commentary below. The doxology cue
\latin{Glória Patri} belongs to the Ordinary rather than to this formulary and
is a scriptureless liturgical formula; no English is supplied for it here.}

\elementguard
\subsubsection*{2. Collect \cue{Coll.}}

\begin{appointedlatin}{\latin{Oratio} --- marginal no.\ 1543}
Omnípotens sempitérne Deus, qui, abundántia pietátis tuæ, et mérita súpplicum
excédis et vota: effúnde super nos misericórdiam tuam; ut dimíttas quæ
consciéntia métuit, et adícias quod orátio non præsúmit. Per Dóminum nostrum.
\end{appointedlatin}

\begin{namedtranslation}{English: Cummiskey (Philadelphia, 1861), Collect of this formulary, printed p.~418}
O Almighty and eternal God, who, in the abundance of thy goodness, exceedest
both the merits and requests of thy suppliants; pour forth thy mercy upon us:
and both pardon what our consciences dreadeth, and grant such blessings as we
dare not presume to ask. Thro'.
\end{namedtranslation}

\englishgap{One shifted subject, declared and left standing. The Latin's last
clause makes the \textit{prayer} the thing that does not presume ---
\latin{quod orátio non præsúmit}, ``what prayer does not presume'' --- and
asks God to \latin{adícias}, to \textit{add} it. The 1861 English makes the
petitioners the ones who dare not ask. The rendering is quoted because it is
the registered public-domain witness; the argument built on the prayer's
grammar is conducted on the Latin, in the element commentary below. The 1861
``dreadeth'' after ``consciences'' is as printed.}

\elementguard
\subsubsection*{3. Epistle \cue{Ep.}}

\begin{appointedlatin}{\latin{Léctio Epístolæ beáti Pauli Apóstoli ad Corínthios} --- 1 Cor. 15, 1-10 --- marginal no.\ 1544}
Fratres: Notum vobis fácio Evangélium, quod prædicávi vobis, quod et
accepístis, in quo et statis, per quod et salvámini: qua ratióne prædicáverim
vobis, si tenétis, nisi frustra credidístis. Trádidi enim vobis in primis,
quod et accépi: quóniam Christus mórtuus est pro peccátis nostris secúndum
Scriptúras: et quia sepúltus est, et quia resurréxit tértia die secúndum
Scriptúras: et quia visus est Cephæ, et post hoc úndecim. Deínde visus est
plus quam quingéntis frátribus simul, ex quibus multi manent usque adhuc,
quidam autem dormiérunt. Deínde visus est Iacóbo, deínde Apóstolis ómnibus:
novíssime autem ómnium tamquam abortívo, visus est et mihi. Ego enim sum
mínimus Apostolórum, qui non sum dignus vocári Apóstolus, quóniam persecútus
sum Ecclésiam Dei. Grátia autem Dei sum id quod sum, et grátia eius in me
vácua non fuit.
\end{appointedlatin}

\begin{namedtranslation}{English: Douay--Rheims (Challoner), 1 Cor. 15:1--10, to the point at which the lesson ends}
Now I make known unto you, brethren, the gospel which I preached to you, which
also you have received and wherein you stand. By which also you are saved, if
you hold fast after what manner I preached unto you, unless you have believed
in vain. For I delivered unto you first of all, which I also received: how
that Christ died for our sins, according to the scriptures: And that he was
buried: and that he rose again the third day, according to the scriptures:
And that he was seen by Cephas, and after that by the eleven. Then was he seen
by more than five hundred brethren at once: of whom many remain until this
present, and some are fallen asleep. After that, he was seen by James: then by
all the apostles. And last of all, he was seen also by me, as by one born out
of due time. For I am the least of the apostles, who am not worthy to be
called an apostle, because I persecuted the church of God. But by the grace of
God, I am what I am. And his grace in me hath not been void \dots
\end{namedtranslation}

\englishgap{Two boundaries, both the missal's. The lesson opens on the
missal's own address: where it prints \latin{Fratres:}, the Clementine reads
\latin{Notum autem vobis fácio, fratres} --- Paul's ``brethren'' stands
mid-clause, and the Douay accordingly begins ``Now I make known unto you,
brethren.'' And the lesson ends mid-verse: the Douay's v.~10 continues ``but I
have laboured more abundantly than all they. Yet not I, but the grace of God
with me,'' and the missal cuts before it, so the reading's last word is grace
not made void, not the Apostle's labour. The same ending stands in the 1570
and 1862 Missals; it is the received lectionary boundary. Neither clause is
supplied here.}

\elementguard
\subsubsection*{4. Gradual \cue{Grad.}}

\begin{appointedlatin}{\latin{Graduale} --- Ps. 27, 7 et 1 --- marginal no.\ 1545}
In Deo sperávit cor meum, et adiútus sum: et reflóruit caro mea, et ex
voluntáte mea confitébor illi. ℣.~Ad te, Dómine, clamávi: Deus meus, ne
síleas, ne discédas a me.
\end{appointedlatin}

\begin{namedtranslation}{English: Douay--Rheims (Challoner), Ps. 27:7 and 27:1 in full, with the sung words in bold}
\textbf{v.~7} The Lord is my helper and my protector: \textbf{in him hath my
heart confided, and I have been helped. And my flesh hath flourished again,
and with my will I will give praise to him.} \textbf{v.~1} A psalm for David
himself. \textbf{Unto thee will I cry, O Lord: O my God, be not thou silent to
me:} lest if thou be silent to me, I become like them that go down into the
pit.
\end{namedtranslation}

\englishgap{Three declarations. \textit{First}, the chant opens \latin{In Deo
sperávit} where the Clementine reads \latin{in ipso sperávit} --- ``in
him,'' resuming the verse's own opening \latin{Dóminus adiútor meus} which the
chant does not sing --- so the Douay's ``in him'' renders a pronoun the chant
replaces with the divine name. \textit{Second}, the versicle sings
\latin{clamávi}, ``I have cried,'' where the Clementine has the future
\latin{clamábo} and the Douay accordingly ``will I cry.'' \textit{Third}, the
versicle's closing \latin{ne discédas a me}, ``depart not from me,'' has no
counterpart in the Clementine verse and therefore none in the Douay; the gap
is declared and left in Latin. One American-edition difference falls inside
the unsung clause printed above: the 1899 printing reads ``lest thou be
silent'' where this e-text reads ``lest if thou be silent.''}

\elementguard
\subsubsection*{5. Alleluia \cue{All.}}

\begin{appointedlatin}{Alleluia, no separate heading --- Ps. 80, 2-3 --- marginal no.\ 1546}
Allelúia, allelúia. ℣. \textnormal{Ps. 80, 2-3} Exsultáte Deo adiutóri
nostro, iubiláte Deo Iacob: súmite psalmum iucúndum cum cíthara. Allelúia.
\end{appointedlatin}

\begin{namedtranslation}{English: Douay--Rheims (Challoner), Ps. 80:2--3 in full, with the sung words in bold}
\textbf{v.~2} \textbf{Rejoice to God our helper: sing aloud to the God of
Jacob.} \textbf{v.~3} \textbf{Take a psalm,} and bring hither the timbrel: the
pleasant psaltery \textbf{with the harp.}
\end{namedtranslation}

\englishgap{The verse condenses Ps. 80:3: the Clementine reads \latin{Súmite
psalmum, et date tympanum; psaltérium iucúndum, cum cíthara}, and the chant
drops the timbrel and the psaltery, attaching \latin{iucúndum} --- ``pleasant''
--- to the psalm where the Bible attaches it to the psaltery. The Douay verse
is therefore printed whole with the sung fragments marked; read in the chant's
own order they say: take a pleasant psalm with the harp. No English join is
composed for it.}

\elementguard
\subsubsection*{6. Gospel \cue{Gosp.}}

\begin{appointedlatin}{\latin{✠ Sequéntia sancti Evangélii secúndum Marcum} --- Marc. 7, 31-37 --- marginal no.\ 1547}
In illo témpore: Exiens Iesus de fínibus Tyri, venit per Sidónem ad mare
Galilǽæ, inter médios fines Decapóleos. Et addúcunt ei surdum et mutum, et
deprecabántur eum, ut impónat illi manum. Et apprehéndens eum de turba
seórsum, misit dígitos suos in aurículas eius: et éxspuens, tétigit linguam
eius: et suspíciens in cælum, ingémuit, et ait illi: Ephphetha, quod est
adaperíre. Et statim apértæ sunt aures eius, et solútum est vínculum linguæ
eius, et loquebátur recte. Et præcépit illis, ne cui dícerent. Quanto autem
eis præcipiébat, tanto magis plus prædicábant: et eo ámplius admirabántur,
dicéntes: Bene ómnia fecit: et surdos fecit audíre, et mutos loqui.
\end{appointedlatin}

\begin{namedtranslation}{English: Douay--Rheims (Challoner), Mk. 7:31--37}
And again going out of the coasts of Tyre, he came by Sidon to the sea of
Galilee, through the midst of the coasts of Decapolis. And they bring to him
one deaf and dumb: and they besought him that he would lay his hand upon him.
And taking him from the multitude apart, he put his fingers into his ears:
and spitting, he touched his tongue. And looking up to heaven, he groaned and
said to him: Ephpheta, which is, Be thou opened. And immediately his ears were
opened and the string of his tongue was loosed and he spoke right. And he
charged them that they should tell no man. But the more he charged them, so
much the more a great deal did they publish it. And so much the more did they
wonder, saying: He hath done all things well. He hath made both the deaf to
hear and the dumb to speak.
\end{namedtranslation}

\englishgap{The missal's \latin{In illo témpore: Exiens Iesus de fínibus
Tyri} names a subject Mark does not name at v.~31 and drops the evangelist's
\latin{íterum} --- the ``again'' that ties this journey back to the earlier
Tyre episode --- which the Douay's opening ``And again going out'' renders.
The liturgical incipit is the missal's own and is not translated here. One
transliteration difference is declared rather than harmonised: the missal
prints \latin{Ephphetha}, with the second aspirate, and the Douay prints
``Ephpheta.''}

\appointedrubric{Credo.}{Printed run-on, right-aligned at the end of the
Gospel's last line. A direction, not a text: the Creed itself is in the
Ordinary.}

\elementguard
\subsubsection*{7. Offertory \cue{Off.}}

\begin{appointedlatin}{\latin{Ant. ad Offertorium} --- Ps. 29, 2-3 --- marginal no.\ 1548}
Exaltábo te, Dómine, quóniam suscepísti me, nec delectásti inimícos meos super
me: Dómine, clamávi ad te, et sanásti me.
\end{appointedlatin}

\begin{namedtranslation}{English: Douay--Rheims (Challoner), Ps. 29:2--3 in full, with the sung words in bold}
\textbf{v.~2} \textbf{I will extol thee, O Lord, for thou hast upheld me: and
hast not made my enemies to rejoice over me.} \textbf{v.~3} \textbf{O Lord} my
God, \textbf{I have cried to thee, and thou hast healed me.}
\end{namedtranslation}

\englishgap{The antiphon sings v.~2 whole and drops \latin{Deus meus} from
v.~3's opening vocative: the Clementine reads \latin{Dómine Deus meus,
clamávi ad te}, and the Douay's ``O Lord my God'' renders the words the chant
does not carry. The verses are printed whole with the sung words marked, and
the omission is declared rather than smoothed over.}

\elementguard
\subsubsection*{8. Secret \cue{Sec.}}

\begin{appointedlatin}{\latin{Secreta} --- marginal no.\ 1549}
Réspice, Dómine, quǽsumus, nostram propítius servitútem: ut, quod offérimus,
sit tibi munus accéptum, et sit nostræ fragilitátis subsídium. Per Dóminum.
\end{appointedlatin}

\begin{namedtranslation}{English: Cummiskey (Philadelphia, 1861), Secret of this formulary, printed p.~420}
Look down, O Lord, we beseech thee, on our homage: that the offerings we make
may be acceptable to thee, and a help to our weakness. Thro'.
\end{namedtranslation}

The 1861 book prints no incipit for this Secret, so its identification was
confirmed clause for clause against the appointed Latin --- ``our homage''
answering \latin{nostram \dots\ servitútem}, ``a help to our weakness''
answering \latin{nostræ fragilitátis subsídium} --- and against the
neighbouring Sundays' orations to prove the filing does not slip by one.

\appointedrubric{Præfatio de Ssma Trinitate.}{Printed immediately after the
Secret, with \latin{Ssma} under a tilde abbreviation for \latin{Sanctissima}:
the Preface of the Most Holy Trinity, which is in the Ordinary. Expanded here
from the same book's own abbreviation practice.}

\elementguard
\subsubsection*{9. Communion \cue{Comm.}}

\begin{appointedlatin}{\latin{Ant. ad Communionem} --- Prov. 3, 9-10 --- marginal no.\ 1550}
Honóra Dóminum de tua substántia, et de primítiis frugum tuárum: et
implebúntur hórrea tua saturitáte, et vino torculária redundábunt.
\end{appointedlatin}

\begin{namedtranslation}{English: Douay--Rheims (Challoner), Prov. 3:9--10 in full, with the sung words in bold}
\textbf{v.~9} \textbf{Honour the Lord with thy substance, and} give him
\textbf{of the first of} all \textbf{thy fruits;} \textbf{v.~10} \textbf{And
thy barns shall be filled with abundance, and thy presses shall run over with
wine.}
\end{namedtranslation}

\englishgap{The antiphon compresses the Clementine at three points: v.~9
loses \latin{ómnium} (``all'') and the closing imperative \latin{da ei}
(``give him''), and v.~10 loses the second \latin{tua}. The Douay verses are
printed whole with the sung words marked; the antiphon's own reading is the
Latin above, and no English join is composed for it.}

\elementguard
\subsubsection*{10. Postcommunion \cue{Postcomm.}}

\begin{appointedlatin}{\latin{Postcommunio} --- marginal no.\ 1551}
Sentiámus, quǽsumus, Dómine, tui perceptióne sacraménti, subsídium mentis et
córporis: ut, in utróque salváti, cæléstis remédii plenitúdine gloriémur. Per
Dóminum.
\end{appointedlatin}

\begin{namedtranslation}{English: Cummiskey (Philadelphia, 1861), Postcommunion of this formulary, printed p.~420}
May we receive, by the participation of these thy mysteries, we beseech thee,
O Lord, help in body and mind: that in the salvation of both we may enjoy the
full effect of this heavenly remedy. Thro'.
\end{namedtranslation}

\englishgap{Two shifts, declared and left standing. The Latin's pair is
\latin{mentis et córporis} --- mind first, body second --- and the 1861
English inverts it to ``body and mind.'' And the Latin's last verb is
\latin{gloriémur}, ``may we glory'' in the fullness of the heavenly remedy;
the 1861 ``enjoy the full effect'' keeps the petition and loses the boast.
The prayer's grammar is discussed on the Latin in the element commentary below.}

\clearpage
